%! Author = Alek
%! Date = 5/3/2026

\documentclass{article}
\usepackage{graphicx} % Required for inserting images
\usepackage{amsmath}
\usepackage[margin=2cm]{geometry}
\usepackage{gensymb}
\usepackage{float}

\title{Multi-Variable Control Systems Project-Update 4}
\author{Alek Krupka}
\date{May 2026}

\begin{document}
    \maketitle

    \section{Introduction}\label{sec:introduction}

    In this project update, system identification methods will be used to perform a frequency
    analysis on data from a system, fit the data to a model, compare the identified model to the
    real model.
    Noise is then added to the model, and the system is re-evaluated to see noise's effect on
    the identified system.
    Finally, a robust controller will be designed that includes the model identification error.

    \section{Plant Model}\label{sec:plant_model}

    The plant model was implemented using the equations derived in previous project updates
    but this time through a Simulink model.
    Using a Simulink model is useful for the following two reasons.

    \begin{enumerate}
        \item It allows for easily adjusting the parameters of the plant through the plant mask as shown by Figure~\ref{fig:open_loop_plant_measurements}.
        \item It allows for easily performing a frequency analysis to create a frd model.
    \end{enumerate}

    First, the modelled plant is shown below by

    \begin{figure}[H]
        \centering
        \includegraphics[width=\linewidth]{ProjectUpdate4Figures/drone_system_model_simulink}
        \caption{Drone Simulink Model}
        \label{fig:drone_system_model_simulink}
    \end{figure}

    And the open loop system containing the plant is given by Figure~\ref{fig:open_loop_plant_measurements}.
    As observed, the parameters are given to the right and can easily be adjusted.

    \begin{figure}[H]
        \centering
        \includegraphics[width=\linewidth]{C:/Users/Alek/Documents/MultivariableControlSystemsProject/ProjectUpdate4/ProjectUpdate4Figures/open_loop_plant_measurements}
        \caption{High Level Simulink Model with Parameter Options}
        \label{fig:open_loop_plant_measurements}
    \end{figure}

    A lot of the work can be done through Simulink's GUI which will be useful for the upcoming analysis.

    \section{Finding the Frequency Response Data}\label{sec:freq_resp_data}

    Using the model, a created a chirp signal can now be given to the plant to create
    data for the system which will then be used for system identification.
    Simulink has a very convenient way of collecting this data through the "frequency response estimator"
    app.
    Performing a frequency analysis using the chirp signal can be done through the following steps in Simulink.

    \begin{enumerate}
        \item Go to the ESTIMATOR TAB.
        \item Select "Input Signal" in the top.
        \item Select "Chirp" from the dropdown
    \end{enumerate}

    Selecting "Chirp" allows for the frequency response to be estimated using the chirp frequency.
    Note that other options can be used.
    However, the chirp input was selected since it covers a wide range and frequencies and since it is
    a simulation, we don't really care about the long time it takes to analyze using this method.

    The chirp signal additionally allows for the model to use different parameters for the frequency
    sweep as shown in Figure~\ref{fig:chirp_signal_frequency_sweep_parameters}

    \begin{figure}[H]
        \centering
        \includegraphics[,width=0.5\linewidth]{ProjectUpdate4Figures/chirp_signal_frequency_sweep_parameters}
        \caption{Frequency Sweep Selected Parameters}
        \label{fig:chirp_signal_frequency_sweep_parameters}
    \end{figure}

    As shown, the frequency is swept between 1Hz and 1kHz.
    While higher frequencies could be analyzed, it is more realistic to keep the sampling time
    relatively low since the sensors are most likely very slow.
    In reality, when performing a system identification on a model like this, the frequency
    would be restricted by sensor speed and interrupt time within the microcontroller.

    Next, Simulink will automatically generate the frd model from the data collected through
    the chirp input.
    This frd model can then be exported to the MATLAB workspace where further analysis can be performed.
    Note that for the MATLAB live script to function, the \text{"chirp\_frd"} variable must be exported to
    the MATLAB workspace.
    The frd model is essentially data taken at a variety of frequnecies and shows the channels between
    inputs and outputs.
    Once this is done though, the frd model can be converted to a state space using the MATLAB
    system identification toolbox.

    \section{System Estimation}\label{sec:estimation}

    Once the frd model is located in the MATLAB workspace there are two options to obtain the
    estimated state space of the system.

    \begin{enumerate}
        \item Use the built-in MATLAB function ssest.
        \item Use the built-in MATLAB function n4sid.
    \end{enumerate}

    Note that the fitfrd function cannot be used here due to the frd being one of a MIMO control system.

    \subsection{Model Order Estimation}\label{subsec:model_order_est}

    To estimate the model order using these functions, the following approach can be used
    as specified in [1].
    Run the ssest/n4sid function where the second parameter is an array from 1-n where n is the max
    order to test.
    In my case, I tested 200 states.
    States with small Henkel singular values are eliminated, and the final model order only contains
    the relatively large singular values.
    However, MATLAB opens a pop-up window and just recommended a model order to use as shown by
    Figure~\ref{fig:ssest_model_order_estimation}.

    \begin{figure}[H]
        \centering
        \includegraphics[width=\linewidth]{C:/Users/Alek/Documents/MultivariableControlSystemsProject/ProjectUpdate4/ProjectUpdate4Figures/ssest_model_order_estimation}
        \caption{Model Order Estimation using ssest}
        \label{fig:ssest_model_order_estimation}
    \end{figure}

    As shown, the recommended model order is 26 states.
    One problem though is that this might be too high of an order and will cause
    too much noise in the model.
    However, the exact model order will be decided at a later time when developing the controller

    \subsection{The Estimated State Space}\label{subsec:estimated_state_space}

    Now that the command is finally run and the model order is estimated the state space is generated.
    Note that because the model is of relatively high order, the exact space needs to be observed from the MATLAB
    live script attached in the appendix.

    \subsection{Comparing the Estimated System to the True System}\label{subsec:est_to_true_system}

    While comparing the estimated to the actual system is almost impossible, the bode plots of both
    MIMO systems can easily be compared.

    The original system's state space can be generated through Simulink's Linear Analysis tool.
    The tool linearizes around an operating point and automatically generates the bode plots.
    Note that the operating point selected is the initial conditions of the plant where there
    is no velocity and the position is above ground.
    This also calculates the bode plot for every channel of the MIMO system as shown by
    Figure~\ref{fig:linearizedopenloopdronemodel}.

    \begin{figure}[H]
        \centering
        \includegraphics[width=\linewidth]{ProjectUpdate4Figures/LinearizedOpenLoopDroneModel}
        \caption{Bode Plot for the Linearized Open Loop Drone Plant}
        \label{fig:linearizedopenloopdronemodel}
    \end{figure}

    The bode plot can again be generated for the identified state space model as shown
    by Figure~\ref{fig:identified_state_space_model}.

    \begin{figure}[H]
        \centering
        \includegraphics[width=\linewidth]{ProjectUpdate4Figures/identified_state_space_model}
        \caption{Bode Plot for Identified State Space Model}
        \label{fig:identified_state_space_model}
    \end{figure}

    As can be observed, the bode plots from the linearized model (Figure~\ref{fig:linearizedopenloopdronemodel})
    look very different from the bode plots generated by the identified model (Figure~\ref{fig:identified_state_space_model}).
    The result is not surprising though, as the default operating point for the drone is upright, meaning
    the linearization is not super accurate in describing the state space model.
    When identifying the system, the drone deviates from the assumed operating point causing
    the bode plots to look more natural.

    \section{Simulating Noise on the Output}\label{sec:simulating_output_noise}

    Now, we add measurment noise to the output of the system when taking the frd response.
    To do this, white noise was added to each sensor output.
    How the Simulink model was changed to accommodate for adding noise is given by Figure~\ref{fig:simulink_model_adjusted_for_noise}.

    \begin{figure}[H]
        \centering
        \includegraphics[width=\linewidth]{C:/Users/Alek/Documents/MultivariableControlSystemsProject/ProjectUpdate4/ProjectUpdate4Figures/simulink_model_adjusted_for_noise}
        \caption{Simulink Model with White Noise}
        \label{fig:simulink_model_adjusted_for_noise}
    \end{figure}

    The same steps as before we repeated to find the frd of the system with output noise.
    The spectral densities for each band of white noise is given by the table below.\\

    \begin{table}[H]
        \centering
    \begin{tabular}{|c|c|}
        \hline
        Sensor & Spectral Density\\
        \hline
        $\theta$ & 0.1\\
        \hline
        $x_{pos}$ & 3\\
        \hline
        $z_{pos}$ & 5\\
        \hline
        speed & 2\\
        \hline
    \end{tabular}
    \label{tab:noise_density_table}
    \end{table}

    Performing the same steps as below the frd produces the bode plot below.

    \begin{figure}[H]
        \centering
        \includegraphics[width=\linewidth]{ProjectUpdate4Figures/noisy_model_bode_plots}
        \caption{Noisy Model Bode Plots}
        \label{fig:noisy_model_bode_plots}
    \end{figure}

    As observed, the bode plots are fairly similar except the phase margin is far worse than desired.
    This again makes sense though because the system is not inherently stable and is incredibly non-linear.

    Next, the chirp signal is run on the estimated state space to generate a new frd.
    This will be used in the code to design a robust controller.

    \section{Designing a Robust Controller using the frd}\label{sec:robust_design}

    Next, a robust controller is designed to stabilize the system, closing the loop and allowing
    it to be controller by a reference.
    Using the frequency response data the problem can be formulated as an $H\infty$ problem with
    uncertainty where the uncertainty is the frequency response error between the new frd
    and the originally derived frd.
    See the MATLAB script on GITHUB for the final controller design and code for
    developing the robust controller.

    \newpage{}
    \section{Sources}\label{sec:sources}

    \begin{enumerate}
        \item https://www.mathworks.com/help/ident/ref/n4sid.html
    \end{enumerate}

    \newpage{}
    \section{Github Repo}\label{sec:repo}

    https://github.com/NightHydra/Multivariable-Control-Systems-Drone-Controller-Project






    \end{document}